\documentclass[a4paper,12pt]{article}

\usepackage[a4paper,margin=16mm]{geometry}
\usepackage[german,main=french]{babel}
\usepackage{fontspec}
\usepackage{microtype}

\setmainfont{EBGaramond}
[
  Extension      = .otf,
  UprightFont    = *-Regular,
  ItalicFont     = *-Italic,
  BoldFont       = *-Bold,
  BoldItalicFont = *-BoldItalic,
  Numbers        = Lowercase
]

\usepackage{paracol}
\newcommand\chunks[2]{%
\begin{leftcolumn*}\begin{otherlanguage}{german}%
{#1}%
\end{otherlanguage}\end{leftcolumn*}%
\begin{rightcolumn}\begin{otherlanguage}{french}%
{#2}%
\end{otherlanguage}\end{rightcolumn}%
}

\newcommand\strophe[1]{\subsection*{\normalsize\normalfont\itshape #1}}

\begin{document}\thispagestyle{empty}
\setlength{\columnsep}{8mm}


\begin{paracol}{2}

  \chunks{\section*{In diesen heil’gen Hallen}}
         {\section*{Dans ces salles sacrées}}

  \chunks{\noindent\raggedright\small Wolfgang Amadeus Mozart, \textit{Die Zauberflöte}}
         {\noindent\raggedright\small Wolfgang Amadeus Mozart, \textit{La Flûte enchantée}}

  \chunks{\noindent\raggedright\small Libretto von Emanuel Schikaneder (1791)}
         {\noindent\raggedright\small Livret d’Emanuel Schikaneder (1791)}

  \chunks{}
         {\noindent\raggedright\small Texte français : Jean Spenlehauer}

  \chunks{\strophe{Zweiter Aufzug, zwölfter Auftritt, Nr.~15, Arie}}
         {\strophe{Deuxième acte, douzième scène, n\textsuperscript{o}~15, air}}

  \chunks{\begin{verse}
          \textsc{Sarastro}\\[\medskipamount]
          In diesen heil’gen Hallen,\\
          Kennt man die Rache nicht!\\
          Und ist ein Mensch gefallen,\\
          Führt Liebe ihn zur Pflicht.\\
          Dann wandelt er an Freundes Hand,\\
          Vergnügt und froh ins bessre Land.
          \end{verse}}
         {\begin{verse}
          \textsc{Sarastro}\\[\medskipamount]
          Dans ces salles sacrées\\
          On ne connaît pas la vengeance !\\
          Et si un homme a chuté ;\\
          L’amour le guide vers son devoir.\\
          Alors il marche guidé par la main d’un ami\\
          Gai et joyeux vers une terre meilleure.
          \end{verse}}

  \chunks{\begin{verse}
          In diesen heil’gen Mauern,\\
          Wo Mensch den Menschen liebt,\\
          Kann kein Verräter lauern,\\
          Weil man dem Feind vergibt.\\
          Wen solche Lehren nicht erfreun,\\
          Verdienet nicht ein Mensch zu sein.
          \end{verse}}
         {\begin{verse}
          Dans ces murs saints\\
          Où l’homme aime l’homme,\\
          Aucun traître ne peut agir,\\
          Parce qu’on pardonne à l’ennemi.\\
          Qui ne se réjouit pas de telles leçons,\\
          Ne mérite pas d’être un homme.
          \end{verse}}
\end{paracol}
\end{document}
